We now add a graded modality to
Graded Base thereby arriving at \textbf{Graded~Modal~Base}. This captures the type system
features of calculi like \textsc{GraD}~\cite{grad} and a monomorphic
subset of~\citet{DBLP:journals/pacmpl/AbelB20}, and permits a simple translation to Linear Base.

\subsection{Graded Modal Base}\label{sec:graded-modal-base}
%
Graded Modal Base can now share the syntax (and operational semantics) of Linear Base:
%
\begin{align*}
    \SYSTEMnt{t} & ::= \SYSTEMmv{x} \mid  \SYSTEMnt{t_{{\mathrm{1}}}} \,  \SYSTEMnt{t_{{\mathrm{2}}}}  \mid  \lambda  \SYSTEMmv{x}  .  \SYSTEMnt{t}  \mid  \textcolor{coeffectColor}{[}  \SYSTEMnt{t}  \textcolor{coeffectColor}{]}  \mid  \mathsf{let} \, \textcolor{coeffectColor}{[}  \SYSTEMmv{x}  \textcolor{coeffectColor}{]} =  \SYSTEMnt{t_{{\mathrm{1}}}}  \, \mathsf{in} \,  \SYSTEMnt{t_{{\mathrm{2}}}} 
    \tag{shared terms} \\
  \SYSTEMnt{A} & ::=  \SYSTEMnt{A}  \xrightarrow{\textcolor{coeffectColor}{ \SYSTEMnt{r} } }  \SYSTEMnt{B}  \mid  \textcolor{coeffectColor}{\square_{ \SYSTEMnt{r} } }  \SYSTEMnt{A}  \mid  \mathrm{K} 
    \tag{extended types}
\end{align*}
%
%
Typing judgments in Graded Modal Base are of the form $ \Delta  \vdash_{\textsc{g}_\square }  \SYSTEMnt{t}  :  \SYSTEMnt{A} $.
The typing rules are the same as Graded Base with addition of typing for
graded modal introduction and elimination:
%
\begin{align*}
\begin{array}{c}
\SYSTEMdruleGradBoxpr{}
\quad
\SYSTEMdruleGradBoxlet{}
\end{array}
\end{align*}
%
The equational theory extends to include $\beta\eta$-equality on
graded modalities (see Appendix~\refappendix{app:definitions-graded-modal-base}).

\subsection{Translation}

\subsubsection{Linear Base to Graded Modal Base translation}
\label{subsec:transl-lb2gmb}

We can now translate Linear Base to Graded \emph{Modal} Base as
a homomorphism on terms, rather than the CPS encoding of Section~\ref{sec:transl}:
%
\renewcommand{\triangleq}{& \oldTriangleq &}
\renewcommand{\equiv}{& \oldEquiv &}
\renewcommand{\SYSTEMdrule}[4][]{#3 \qquad #2}
%
\begin{gather*}
{\setlength{\arraycolsep}{0.1em}
\begin{array}{lcl@{\hspace{2em}}r}
\SYSTEMdruleTranslLtoGTmvar{} \\
\SYSTEMdruleTranslLtoGTmabs{} \\
\SYSTEMdruleTranslLtoGTmapp{} \\
\SYSTEMdruleTranslLtoGTmpr{} \\
\SYSTEMdruleTranslLtoGTmlet{} & \text{(Term translation)} \\[0.5em]
\SYSTEMdruleTranslLtoGTyfunTy{} \\
\SYSTEMdruleTranslLtoGTyboxTy{} \\
\SYSTEMdruleTranslLtoGTybaseTy{} & \text{(Type translation)} \\[0.5em]
\SYSTEMdruleTranslLtoGGtxtempty{} \\
\SYSTEMdruleTranslLtoGGtxtlin{} \\
\SYSTEMdruleTranslLtoGGtxtgrad{} & \text{(Context translation)}
\end{array}}
\end{gather*}
\let\SYSTEMdrule\oldSYSTEMdrule
\let\equiv\oldEquiv
\let\triangleq\oldTriangleq
%
For types, we translate $ \textcolor{coeffectColor}{\square_{ \SYSTEMnt{r} } }     $ in the straightforward way, directly
to the Graded Modal Base graded modality. To translate
the linear function space we grade the resulting function arrow
in Graded Modal Base at $1 \in \mathcal{R}$ and similarly for linear assumptions.

\begin{example}
  \label{exm:k-combinator}
  The following Linear Base judgment types a $K$-combinator-like expression:
  $$   \SYSTEMmv{x}  :  \SYSTEMnt{B}    \vdash_{\textsc{l} }   \lambda  \SYSTEMmv{y}  .   \mathsf{let} \, \textcolor{coeffectColor}{[}  \SYSTEMmv{z}  \textcolor{coeffectColor}{]} =  \SYSTEMmv{y}  \, \mathsf{in} \,  \SYSTEMmv{x}    :     \textcolor{coeffectColor}{\square_{ \SYSTEMsym{0} } }  \SYSTEMnt{A}    \multimap  \SYSTEMnt{B}  $$
  The parameter $y$ is used by eliminating the graded modality,
  binding to $z$ but then discarding the term. Its translation to Graded Modal Base results
  in an identical term but with translated types:
  %
  $$   \SYSTEMmv{x}  :_{\textcolor{coeffectColor}{ \SYSTEMsym{1} } }   \textcolor{LGcolor}{\llparenthesis} \smidge  \SYSTEMnt{B}  \smidge \textcolor{LGcolor}{\rrparenthesis}     \vdash_{\textsc{g} }   \lambda  \SYSTEMmv{y}  .   \mathsf{let} \, \textcolor{coeffectColor}{[}  \SYSTEMmv{z}  \textcolor{coeffectColor}{]} =  \SYSTEMmv{y}  \, \mathsf{in} \,  \SYSTEMmv{x}    :     \textcolor{coeffectColor}{\square_{ \SYSTEMsym{0} } }   \textcolor{LGcolor}{\llparenthesis} \smidge  \SYSTEMnt{A}  \smidge \textcolor{LGcolor}{\rrparenthesis}     \xrightarrow{\textcolor{coeffectColor}{ \SYSTEMsym{1} } }   \textcolor{LGcolor}{\llparenthesis} \smidge  \SYSTEMnt{B}  \smidge \textcolor{LGcolor}{\rrparenthesis}   $$
  %
\end{example}

\begin{restatable}{thrm}{linToGradTranslation}[Soundness of the Linear Base to Graded Modal Base translation]\label{thrm:linToGradTranslation}
  \begin{itemize}[itemsep=0.25em,topsep=0.25em]
    \item Type preservation: $ \Gamma  \vdash_{\textsc{l} }  \SYSTEMnt{t}  :  \SYSTEMnt{A}  \implies   \textcolor{LGcolor}{\llparenthesis}  \Gamma  \textcolor{LGcolor}{\rrparenthesis}   \vdash_{\textsc{g}_\square }   \textcolor{LGcolor}{\llparenthesis} \smidge  \SYSTEMnt{t}  \smidge \textcolor{LGcolor}{\rrparenthesis}   :   \textcolor{LGcolor}{\llparenthesis} \smidge  \SYSTEMnt{A}  \smidge \textcolor{LGcolor}{\rrparenthesis}  $
    \item Operational correspondence: $ \SYSTEMnt{t}  \rightsquigarrow_{\textsc{l} }  \SYSTEMnt{t'}  \implies   \textcolor{LGcolor}{\llparenthesis} \smidge  \SYSTEMnt{t}  \smidge \textcolor{LGcolor}{\rrparenthesis}   \rightsquigarrow_{\textsc{g}\Box }^\ast   \textcolor{LGcolor}{\llparenthesis} \smidge  \SYSTEMnt{t'}  \smidge \textcolor{LGcolor}{\rrparenthesis}  $
    \item Equation preservation: $ \SYSTEMnt{t}  \,\equiv_{\textsc{l} }\,  \SYSTEMnt{t'}  \implies   \textcolor{LGcolor}{\llparenthesis} \smidge  \SYSTEMnt{t}  \smidge \textcolor{LGcolor}{\rrparenthesis}   \,\equiv_{\textsc{g}_\Box}\,   \textcolor{LGcolor}{\llparenthesis} \smidge  \SYSTEMnt{t'}  \smidge \textcolor{LGcolor}{\rrparenthesis}  $
  \end{itemize}
\end{restatable}

\noindent
Appendix~\refappendix{app:proofs-lin-to-grad-mod} provides the proof along with its supporting lemmas.

\begin{remark}
  We might consider whether there is a dual to the above, a kind of completeness result,
i.e., that
given $ \Delta  \vdash_{\textsc{g} }  \SYSTEMnt{t}  :  \SYSTEMnt{A} $
there exists $\Gamma, \SYSTEMnt{t'}, \SYSTEMnt{A'}$
such that $ \Gamma  \vdash_{\textsc{l} }  \SYSTEMnt{t'}  :  \SYSTEMnt{A'} $ and $  \textcolor{LGcolor}{\llparenthesis}  \Gamma  \textcolor{LGcolor}{\rrparenthesis}   \vdash_{\textsc{g} }   \textcolor{LGcolor}{\llparenthesis} \smidge  \SYSTEMnt{t'}  \smidge \textcolor{LGcolor}{\rrparenthesis}   :   \textcolor{LGcolor}{\llparenthesis} \smidge  \SYSTEMnt{A'}  \smidge \textcolor{LGcolor}{\rrparenthesis}  $
and $\Delta \equiv  \textcolor{LGcolor}{\llparenthesis}  \Gamma  \textcolor{LGcolor}{\rrparenthesis} $
and $\SYSTEMnt{t} \equiv  \textcolor{LGcolor}{\llparenthesis} \smidge  \SYSTEMnt{t'}  \smidge \textcolor{LGcolor}{\rrparenthesis} $
and $\SYSTEMnt{A} \equiv  \textcolor{LGcolor}{\llparenthesis} \smidge  \SYSTEMnt{A'}  \smidge \textcolor{LGcolor}{\rrparenthesis} $.
However, an easy counterexample is the following judgment:
$$   \SYSTEMmv{x}  :_{\textcolor{coeffectColor}{ \SYSTEMsym{1} } }  \SYSTEMnt{B}    \vdash_{\textsc{g} }   \lambda  \SYSTEMmv{y}  .  \SYSTEMmv{x}   :   \SYSTEMnt{A}  \xrightarrow{\textcolor{coeffectColor}{ \SYSTEMsym{0} } }  \SYSTEMnt{B}  .$$
There is no type in Linear Base which translates to $ \SYSTEMnt{A}  \xrightarrow{\textcolor{coeffectColor}{ \SYSTEMsym{0} } }  \SYSTEMnt{B} $ as we always
translate the linear function type to $ \SYSTEMnt{A}  \xrightarrow{\textcolor{coeffectColor}{ \SYSTEMsym{1} } }  \SYSTEMnt{B} $. Indeed, the analogous
term in Linear Base was already shown in Example~\ref{exm:k-combinator}, where
we see that the result of the translation has type $   \textcolor{coeffectColor}{\square_{ \SYSTEMsym{0} } }   \textcolor{LGcolor}{\llparenthesis} \smidge  \SYSTEMnt{A}  \smidge \textcolor{LGcolor}{\rrparenthesis}     \xrightarrow{\textcolor{coeffectColor}{ \SYSTEMsym{1} } }   \textcolor{LGcolor}{\llparenthesis} \smidge  \SYSTEMnt{B}  \smidge \textcolor{LGcolor}{\rrparenthesis}  $.

Thus, the translation from Linear Base to Graded Modal Base is not surjective. Neither
is it injective: both the linear assumption $ \SYSTEMmv{x}  :  \SYSTEMnt{A} $ and the graded assumption
$ \SYSTEMmv{x}  : \textcolor{coeffectColor}{[}  \SYSTEMnt{A} {\textcolor{coeffectColor}{]_{ \SYSTEMsym{1} } } } $ are translated to $ \SYSTEMmv{x}  :_{\textcolor{coeffectColor}{ \SYSTEMsym{1} } }   \textcolor{LGcolor}{\llparenthesis} \smidge  \SYSTEMnt{A}  \smidge \textcolor{LGcolor}{\rrparenthesis}  $. For those semirings
in which $1$ does not capture linearity (e.g., the singleton semiring, any semiring
where $0 \leq 1$, and others), linear and non-linear
use is conflated by this translation.
\end{remark}

\subsubsection{A Linearity-Preserving Linear Base to Graded Modal Base Translation}
\label{subsec:linearity-preserving}

Instead, we can give an alternate translation which is injective
and \emph{linearity preserving}: keeping the translation of terms
as above, we can translate types and grades from Linear Base with semiring $\mathcal{R}$
into types and grades in Graded Modal Base with semiring $(\mathcal{R} \times \{\SYSTEMsym{0},
\SYSTEMsym{1},  \omega \})$, i.e., the product semiring (Example~\ref{exm:product})
of $\mathcal{R}$ with the linearity-capturing semiring (Example~\ref{exm:none-one-tons}).
%
\renewcommand{\triangleq}{& \oldTriangleq &}
\renewcommand{\equiv}{& \oldEquiv &}
\renewcommand{\SYSTEMdrule}[4][]{#3 \qquad #2}
%
\begin{gather*}
{\setlength{\arraycolsep}{0.1em}
\begin{array}{lcl@{\hspace{2em}}r}
\SYSTEMdruleTranslLtoGAltTyfunTy{} \\
\SYSTEMdruleTranslLtoGAltTyboxTy{} \\
\SYSTEMdruleTranslLtoGAltTybaseTy{} & \text{(Type translation)} \\[1em]
%\SYSTEMdruleTranslLtoGAltGtxtempty{} \\
\SYSTEMdruleTranslLtoGAltGtxtlin{} \\
\SYSTEMdruleTranslLtoGAltGtxtgrad{} & \text{(Context translation)}
\end{array}}
\end{gather*}
\let\SYSTEMdrule\oldSYSTEMdrule
\let\equiv\oldEquiv
\let\triangleq\oldTriangleq
%
%\newpage
Thus, the translation is now injective in the types since we translate a linear assumption $ \SYSTEMmv{x}  :  \SYSTEMnt{A} $ to
have grade $ \langle  \SYSTEMsym{1}  ,  \SYSTEMsym{1}  \rangle $ but a graded assumption $ \SYSTEMmv{x}  : \textcolor{coeffectColor}{[}  \SYSTEMnt{A} {\textcolor{coeffectColor}{]_{ \SYSTEMnt{r} } } } $ to $ \langle  \SYSTEMnt{r}  ,   \omega   \rangle $
which is a different grade even if $r = 1$. By taking the pairing with the linearity semiring,
we can capture the combination of linearity and grading that occurs in Linear Base, but
now in the Graded Modal Base setting.

  Operational correspondence and equation preservation hold as before in Theorem~\ref{thrm:linToGradTranslation} since the term translation is unchanged. We then re-restablish
  type preservation of the translation:
%
\begin{restatable}{thrm}{linToGradAltTranslation}[Soundness of the Linearity-Preserving Linear Base to Graded Modal Base translation]
  \begin{itemize}[itemsep=0.2em]
    \item Type preservation: $ \Gamma  \vdash_{\textsc{l} }  \SYSTEMnt{t}  :  \SYSTEMnt{A}  \implies   \textcolor{LGcolor}{\llparenthesis}  \Gamma  \textcolor{LGcolor}{\rrparenthesis}'   \vdash_{\textsc{g}_\square }   \textcolor{LGcolor}{\llparenthesis} \smidge  \SYSTEMnt{t}  \smidge \textcolor{LGcolor}{\rrparenthesis}   :   \textcolor{LGcolor}{\llparenthesis} \smidge  \SYSTEMnt{A}  \smidge \textcolor{LGcolor}{\rrparenthesis}'  $
  \end{itemize}
Appendix~\refappendix{app:proofs-lin-to-grad-mod-alt} provides the proof.
\end{restatable}

The key insight in this translation is that Linear Base essentially comprises two parallel
forms of analysis: a graded analysis and a latent linear analysis witnessed by the
corresponding two forms of typing assumption in Linear Base contexts.
Our translation here captures these two aspects systematically,
by taking the product of the same grading semiring and the linearity semiring.

A question may arise here:
\textit{What is the meaning of a grade such as $ \langle  \SYSTEMsym{2}  ,  \SYSTEMsym{1}  \rangle $,
where the first component looks non-linear in $\mathcal{R}$
but the second represents linearity?} Such a grade is not
in the image of our translation but we may still consider what it is
like to type terms in Graded (Modal) Base with $\mathcal{R} \times \{0, 1,
\omega\}$. First, let us decompose $ \langle  \SYSTEMsym{2}  ,  \SYSTEMsym{1}  \rangle $
into $ \langle  \SYSTEMsym{1}  \SYSTEMsym{+}  \SYSTEMsym{1}  ,  \SYSTEMsym{1}  \rangle $
for clarity, i.e., for any semiring $\mathcal{R}$ in the first component.
Such grades could be formed via approximation if
$\mathcal{R}$ permits $ \SYSTEMsym{1}  \, \textcolor{coeffectColor}{\sqsubseteq} \,  \SYSTEMsym{1}  \SYSTEMsym{+}  \SYSTEMsym{1} $ and thus $  \langle  \SYSTEMsym{1}  ,  \SYSTEMsym{1}  \rangle   \, \textcolor{coeffectColor}{\sqsubseteq} \,   \langle  \SYSTEMsym{1}  \SYSTEMsym{+}  \SYSTEMsym{1}  ,  \SYSTEMsym{1}  \rangle  $
is an over-approximation of grades, e.g., useful for when calculating the grade
across multiple program branches. Without such an
approximation, $ \langle  \SYSTEMsym{1}  \SYSTEMsym{+}  \SYSTEMsym{1}  ,  \SYSTEMsym{1}  \rangle $ cannot be derived from variable use
and contraction of variables alone since in $(\mathcal{R} \times \{\SYSTEMsym{0},
\SYSTEMsym{1},  \omega \})$ we have
$ \langle  \SYSTEMsym{1}  ,  \SYSTEMsym{1}  \rangle   \SYSTEMsym{+}   \langle  \SYSTEMsym{1}  ,  \SYSTEMsym{1}  \rangle  =  \langle  \SYSTEMsym{1}  \SYSTEMsym{+}  \SYSTEMsym{1}  ,   \omega   \rangle $.
Such a grade could only arise from an open term applying
a function, e.g.,
$$    \SYSTEMmv{f}  :_{\textcolor{coeffectColor}{ \SYSTEMsym{1} } }    \SYSTEMnt{A}  \xrightarrow{\textcolor{coeffectColor}{  \langle  \SYSTEMsym{1}  \SYSTEMsym{+}  \SYSTEMsym{1}  ,  \SYSTEMsym{1}  \rangle  } }  \SYSTEMnt{B}     ,   \SYSTEMmv{x}  :_{\textcolor{coeffectColor}{  \langle  \SYSTEMsym{1}  \SYSTEMsym{+}  \SYSTEMsym{1}  ,  \SYSTEMsym{1}  \rangle  } }  \SYSTEMnt{A}    \vdash_{\textsc{g} }   \SYSTEMmv{f} \,  \SYSTEMmv{x}   :  \SYSTEMnt{B} .$$
However there is no well-typed closed term that can be formed to substitute into
$f$ since we cannot form $ \langle  \SYSTEMsym{1}  \SYSTEMsym{+}  \SYSTEMsym{1}  ,  \SYSTEMsym{1}  \rangle $ without use of such an open term.
Consequently, for closed terms, we cannot form $ \langle  \SYSTEMsym{1}  \SYSTEMsym{+}  \SYSTEMsym{1}  ,  \SYSTEMsym{1}  \rangle $
unless it is permitted by approximation in $\mathcal{R}$.
In the approximation case, it is up to the choice (or design) of semiring
to permit approximation via the definition of $     \, \textcolor{coeffectColor}{\sqsubseteq} \,     $.

More generally, we may consider the advantage of pairing
semirings with the linearity semiring. Consider a semiring such as the security
lattice semiring (Example~\ref{exm:security-lattice}) where $1 =  \mathsf{Lo} $.
Since $ \mathsf{Lo}  +  \mathsf{Lo}  =  \mathsf{Lo} $ and $ \SYSTEMsym{0}  \, \textcolor{coeffectColor}{\sqsubseteq} \,   \mathsf{Lo}  $ then a $ \mathsf{Lo} $-graded variable
can represent both linear or non-linear usage. The pairing
$ \langle   \mathsf{Lo}   ,  \SYSTEMsym{1}  \rangle $ then distinguishes a linear low-security variable
from $ \langle   \mathsf{Lo}   ,   \omega   \rangle $ which instead represents a non-linear low-security variable.
Conversely, recall that $0 =  \mathsf{Hi} $,
giving the translation $  \textcolor{LGcolor}{\llparenthesis} \smidge   \textcolor{coeffectColor}{\square_{  \mathsf{Hi}  } }  \SYSTEMnt{A}   \smidge \textcolor{LGcolor}{\rrparenthesis}'   \equiv   \textcolor{coeffectColor}{\square_{  \langle   \mathsf{Hi}   ,   \omega   \rangle  } }  \SYSTEMnt{A}  $,
which allows unlimited usage, but limited to $ \mathsf{Hi} $ security contexts.

Our approach here is not necessarily the only sound and injective translation;
however, it seems particularly illuminating,
in that the product construction clearly demarcates and represents the two analyses present within Linear Base:
the explicit coeffect in the grades on the one hand and the latent linearity coeffect of the
system (not present in Graded Base) on the other.

% is it really a necessary construction? No there might be other choices
% Lin-completion is not satisfactory as it can give Lin <= 1 + 1
%

% Reviewer:
% The goal of translation seems to be to /explain/ one world in terms of another, but the graded world doesn't use semi-rings like this ... so is it an explanation or a proof tool that "just happens" to work?

% Our insight: that the Linear Base is really like having a pair of two semirings
% one for R and one for linearity, and this conceptually shows that there are two
% separate kinds of tracking going on , which we then naturally present as a product
% of two semirings in this translation to Graded Base

% Lin <= 1

% <1, 1>

%Does it matter that it seems to make no sense to "translate back" in the other direction?

\subsubsection{Graded Modal Base to Linear Base Translation}
\label{subsec:transl-gmb2lb}

\renewcommand{\triangleq}{& \oldTriangleq &}
\renewcommand{\equiv}{& \oldEquiv &}
\renewcommand{\SYSTEMdrule}[4][]{#3 \qquad #2}

We can now extend the previous translation from Graded Base to Linear Base
(Section~\ref{sec:transl}) to translate from Graded Modal Base to Linear Base:
%
\begin{gather*}
{\setlength{\arraycolsep}{0.1em}
\begin{array}{lcl@{\hspace{2em}}r}
\SYSTEMdruleTranslGtoLTmvar{} \\
\SYSTEMdruleTranslGtoLTmabs{}
\SYSTEMdruleTranslGtoLTmapp{} \\
\SYSTEMdruleTranslGtoLTmpr{} \\
\SYSTEMdruleTranslGtoLTmlet{} & \text{(Term translation)} \\[0.75em]
\SYSTEMdruleTranslGtoLTyfunTy{} \\
\SYSTEMdruleTranslGtoLTyboxTy{} \\
\SYSTEMdruleTranslGtoLTybaseTy{} & \text{(Type translation)} \\[0.75em]
\SYSTEMdruleTranslGtoLGtxtempty{} \\
\SYSTEMdruleTranslGtoLGtxtgrad{} & \text{(Context translation)}
\end{array}}
\end{gather*}
\let\SYSTEMdrule\oldSYSTEMdrule
\let\equiv\oldEquiv
\let\triangleq\oldTriangleq
%
The only new parts are the translation of the graded modality which maps directly
onto the graded modality in Linear Base. We then (re)establish soundness. Appendix~\refappendix{app:proofs-grad-mod-to-lin} provides the proofs.
%
\begin{restatable}{thrm}{gradModToLinTranslation}[Soundness of the Graded Modal Base to Linear Base translation]
\begin{itemize}[itemsep=0.2em]
  \item Type preservation: $ \Delta  \vdash_{\textsc{g}_\square }  \SYSTEMnt{t}  :  \SYSTEMnt{A}  \implies   \textcolor{GLcolor}{\llbracket}  \Delta  \textcolor{GLcolor}{\rrbracket}   \vdash_{\textsc{l} }   \textcolor{GLcolor}{\llbracket}  \SYSTEMnt{t}  \textcolor{GLcolor}{\rrbracket}   :   \textcolor{GLcolor}{\llbracket}  \SYSTEMnt{A}  \textcolor{GLcolor}{\rrbracket}  $
  \item Operational correspondence: $ \SYSTEMnt{t}  \rightsquigarrow_{\textsc{g}\Box}  \SYSTEMnt{t'}  \implies   \textcolor{GLcolor}{\llbracket}  \SYSTEMnt{t}  \textcolor{GLcolor}{\rrbracket}   \rightsquigarrow_{\textsc{l} }^\ast   \textcolor{GLcolor}{\llbracket}  \SYSTEMnt{t'}  \textcolor{GLcolor}{\rrbracket}  $
  \item Equation preservation: $ \SYSTEMnt{t}  \,\equiv_{\textsc{g}_\Box}\,  \SYSTEMnt{t'}  \implies   \textcolor{GLcolor}{\llbracket}  \SYSTEMnt{t}  \textcolor{GLcolor}{\rrbracket}   \,\equiv_{\textsc{l} }\,   \textcolor{GLcolor}{\llbracket}  \SYSTEMnt{t'}  \textcolor{GLcolor}{\rrbracket}  $.
 \end{itemize}
  \label{ref:gradModToLinTranslation}
\end{restatable}
%

\subsubsection{Mutual relationship}
\label{subsec:mutual-relationship}

Given the pair of translations here, we might wonder: what is the relationship between
the translations? Are they mutually inverse? By considering just the translation of function
types we can quickly see that the translations are not inverse:
%
\begin{align*}
 \textcolor{GLcolor}{\llbracket}   \textcolor{LGcolor}{\llparenthesis} \smidge   \SYSTEMnt{A}  \multimap  \SYSTEMnt{B}   \smidge \textcolor{LGcolor}{\rrparenthesis}   \textcolor{GLcolor}{\rrbracket}  & =   \textcolor{GLcolor}{\llbracket}   \SYSTEMnt{A}  \xrightarrow{\textcolor{coeffectColor}{ \SYSTEMsym{1} } }  \SYSTEMnt{B}   \textcolor{GLcolor}{\rrbracket}   \equiv    \textcolor{coeffectColor}{\square_{ \SYSTEMsym{1} } }    \textcolor{GLcolor}{\llbracket}   \textcolor{LGcolor}{\llparenthesis} \smidge  \SYSTEMnt{A}  \smidge \textcolor{LGcolor}{\rrparenthesis}   \textcolor{GLcolor}{\rrbracket}     \multimap    \textcolor{GLcolor}{\llbracket}   \textcolor{LGcolor}{\llparenthesis} \smidge  \SYSTEMnt{B}  \smidge \textcolor{LGcolor}{\rrparenthesis}   \textcolor{GLcolor}{\rrbracket}    
\\
 \textcolor{LGcolor}{\llparenthesis} \smidge   \textcolor{GLcolor}{\llbracket}   \SYSTEMnt{A}  \xrightarrow{\textcolor{coeffectColor}{ \SYSTEMnt{r} } }  \SYSTEMnt{B}   \textcolor{GLcolor}{\rrbracket}   \smidge \textcolor{LGcolor}{\rrparenthesis}  & =   \textcolor{LGcolor}{\llparenthesis} \smidge    \textcolor{coeffectColor}{\square_{ \SYSTEMnt{r} } }  \SYSTEMnt{A}   \multimap  \SYSTEMnt{B}   \smidge \textcolor{LGcolor}{\rrparenthesis}   \equiv    \textcolor{coeffectColor}{\square_{ \SYSTEMnt{r} } }   \textcolor{LGcolor}{\llparenthesis} \smidge   \textcolor{GLcolor}{\llbracket}  \SYSTEMnt{A}  \textcolor{GLcolor}{\rrbracket}   \smidge \textcolor{LGcolor}{\rrparenthesis}    \xrightarrow{\textcolor{coeffectColor}{ \SYSTEMsym{1} } }   \textcolor{LGcolor}{\llparenthesis} \smidge   \textcolor{GLcolor}{\llbracket}  \SYSTEMnt{B}  \textcolor{GLcolor}{\rrbracket}   \smidge \textcolor{LGcolor}{\rrparenthesis}   
\end{align*}
%
Similarly, given $   \SYSTEMmv{x}  :  \SYSTEMnt{A}    \vdash_{\textsc{l} }  \SYSTEMmv{x}  :  \SYSTEMnt{A} $ in Linear Base
then its translation to Graded Modal Base is $   \SYSTEMmv{x}  :_{\textcolor{coeffectColor}{ \SYSTEMsym{1} } }  \SYSTEMnt{A}    \vdash_{\textsc{g}_\square }  \SYSTEMmv{x}  :  \SYSTEMnt{A} $
which is translated back to Linear Base as $   \SYSTEMmv{x}  : \textcolor{coeffectColor}{[}  \SYSTEMnt{A} {\textcolor{coeffectColor}{]_{ \SYSTEMsym{1} } } }    \vdash_{\textsc{l} }  \SYSTEMmv{x}  :  \SYSTEMnt{A} $.

However, by composing operational correspondence theorems, then round-tripping the translations
yields a term that corresponds operationally to the original, in both directions:
%
\begin{corollary}[Observational isomorphism]
If $ \SYSTEMnt{t}  \rightsquigarrow_{\textsc{l} }  \SYSTEMnt{t'} $ then $  \textcolor{GLcolor}{\llbracket}   \textcolor{LGcolor}{\llparenthesis} \smidge  \SYSTEMnt{t}  \smidge \textcolor{LGcolor}{\rrparenthesis}   \textcolor{GLcolor}{\rrbracket}   \rightsquigarrow_{\textsc{l} }^\ast   \textcolor{GLcolor}{\llbracket}   \textcolor{LGcolor}{\llparenthesis} \smidge  \SYSTEMnt{t'}  \smidge \textcolor{LGcolor}{\rrparenthesis}   \textcolor{GLcolor}{\rrbracket}  $
and if $ \SYSTEMnt{t}  \rightsquigarrow_{\textsc{g}\Box}  \SYSTEMnt{t'} $ then $  \textcolor{LGcolor}{\llparenthesis} \smidge   \textcolor{GLcolor}{\llbracket}  \SYSTEMnt{t}  \textcolor{GLcolor}{\rrbracket}   \smidge \textcolor{LGcolor}{\rrparenthesis}   \rightsquigarrow_{\textsc{g}\Box }^\ast   \textcolor{LGcolor}{\llparenthesis} \smidge   \textcolor{GLcolor}{\llbracket}  \SYSTEMnt{t'}  \textcolor{GLcolor}{\rrbracket}   \smidge \textcolor{LGcolor}{\rrparenthesis}  $.
\end{corollary}
%
We might also wonder if there is an adjoint relationship between the translations. If $ \textcolor{GLcolor}{\llbracket}   -   \textcolor{GLcolor}{\rrbracket}  \dashv  \textcolor{LGcolor}{\llparenthesis} \smidge   -   \smidge \textcolor{LGcolor}{\rrparenthesis} $ then we would expect to be able to construct a mapping in the category of types (whose morphisms are derivations of judgments with source as inputs and the target as the result):
%
\[
\eta_A : A \rightarrow  \textcolor{GLcolor}{\llbracket}   \textcolor{LGcolor}{\llparenthesis} \smidge  \SYSTEMnt{A}  \smidge \textcolor{LGcolor}{\rrparenthesis}   \textcolor{GLcolor}{\rrbracket} 
\]
%
However, there is no such mapping for $ \SYSTEMnt{A}  \multimap  \SYSTEMnt{B} $ to $  \textcolor{GLcolor}{\llbracket}   \textcolor{LGcolor}{\llparenthesis} \smidge   \SYSTEMnt{A}  \multimap  \SYSTEMnt{B}   \smidge \textcolor{LGcolor}{\rrparenthesis}   \textcolor{GLcolor}{\rrbracket}   \equiv    \textcolor{coeffectColor}{\square_{ \SYSTEMsym{1} } }   \textcolor{GLcolor}{\llbracket}   \textcolor{LGcolor}{\llparenthesis} \smidge  \SYSTEMnt{A}  \smidge \textcolor{LGcolor}{\rrparenthesis}   \textcolor{GLcolor}{\rrbracket}    \multimap   \textcolor{GLcolor}{\llbracket}   \textcolor{LGcolor}{\llparenthesis} \smidge  \SYSTEMnt{B}  \smidge \textcolor{LGcolor}{\rrparenthesis}   \textcolor{GLcolor}{\rrbracket}   $ in general since this is undefined for a higher-order type.

Consider $  (    \mathrm{K}   \multimap   \mathrm{K}    )   \multimap   \mathrm{K}  $ for the base type $ \mathrm{K} $ whose translation between
the systems is the identity, and thus $ \textcolor{GLcolor}{\llbracket}   \textcolor{LGcolor}{\llparenthesis} \smidge   \mathrm{K}   \smidge \textcolor{LGcolor}{\rrparenthesis}   \textcolor{GLcolor}{\rrbracket}  =  \mathrm{K} $. Then $\eta$ would need to provide
a derivation of the judgment:
%
\begin{align*}
  (    \mathrm{K}   \multimap   \mathrm{K}    )   \multimap   \mathrm{K}   \vdash   \textcolor{coeffectColor}{\square_{ \SYSTEMsym{1} } }   (    \textcolor{coeffectColor}{\square_{ \SYSTEMsym{1} } }   \mathrm{K}    \multimap   \mathrm{K}    )    \multimap   \mathrm{K}  
\end{align*}
%
Such a derivation would require us to map a value of type $ \textcolor{coeffectColor}{\square_{ \SYSTEMsym{1} } }   (    \textcolor{coeffectColor}{\square_{ \SYSTEMsym{1} } }   \mathrm{K}    \multimap   \mathrm{K}    )  $
to one of type $  \mathrm{K}   \multimap   \mathrm{K}  $. We can certainly derive $ \textcolor{coeffectColor}{\square_{ \SYSTEMsym{1} } }   (    \textcolor{coeffectColor}{\square_{ \SYSTEMsym{1} } }   \mathrm{K}    \multimap   \mathrm{K}    )   \vdash   \textcolor{coeffectColor}{\square_{ \SYSTEMsym{1} } }   \mathrm{K}    \multimap   \mathrm{K}  $ by elimination
and dereliction, but $  \textcolor{coeffectColor}{\square_{ \SYSTEMsym{1} } }   \mathrm{K}    \multimap   \mathrm{K}   \not\vdash   \mathrm{K}   \multimap   \mathrm{K}  $. To see why, consider
the single-point semiring $0 = 1$ and thus a function typed $  \textcolor{coeffectColor}{\square_{ \SYSTEMsym{1} } }   \mathrm{K}    \multimap   \mathrm{K}  $ may not use its argument
at all, whereas $  \mathrm{K}   \multimap   \mathrm{K}  $ must use its argument exactly once.

Given the above negative result, we could then consider the opposite adjoint relationship, i.e., perhaps
 $ \textcolor{LGcolor}{\llparenthesis} \smidge   -   \smidge \textcolor{LGcolor}{\rrparenthesis}  \dashv  \textcolor{GLcolor}{\llbracket}   -   \textcolor{GLcolor}{\rrbracket} $? In which case, we would require the counit
 operation:
\[
\epsilon_A :  \textcolor{GLcolor}{\llbracket}   \textcolor{LGcolor}{\llparenthesis} \smidge  \SYSTEMnt{A}  \smidge \textcolor{LGcolor}{\rrparenthesis}   \textcolor{GLcolor}{\rrbracket}  \to A
\]
%
However, there is no such mapping, even in a first-order case: for
$  \textcolor{GLcolor}{\llbracket}   \textcolor{LGcolor}{\llparenthesis} \smidge    \mathrm{K}   \multimap   \mathrm{K}    \smidge \textcolor{LGcolor}{\rrparenthesis}   \textcolor{GLcolor}{\rrbracket}   \equiv    \textcolor{coeffectColor}{\square_{ \SYSTEMsym{1} } }   \textcolor{GLcolor}{\llbracket}   \textcolor{LGcolor}{\llparenthesis} \smidge   \mathrm{K}   \smidge \textcolor{LGcolor}{\rrparenthesis}   \textcolor{GLcolor}{\rrbracket}    \multimap   \textcolor{GLcolor}{\llbracket}   \textcolor{LGcolor}{\llparenthesis} \smidge   \mathrm{K}   \smidge \textcolor{LGcolor}{\rrparenthesis}   \textcolor{GLcolor}{\rrbracket}   $ we need to
map to $  \mathrm{K}   \multimap   \mathrm{K}  $ but as described above, no such derivation is possible.

Thus the translations are neither inverses nor adjoint. We leave it as future work to
consider whether more complex invertible translations are possible.

\subsection{Generalised $\Box$-elimination}
\label{sec:flattening}
Elimination of graded modalities in both linear- and graded-base systems looks almost identical,
eliminating a graded modality at grade $r$ by a kind of
`cut' (substitution) against a variable with matching grade $r$:
%
\begin{gather*}
\SYSTEMdruleLinlet{} \quad \SYSTEMdruleGradBoxlet{}
\end{gather*}
%
Some graded-base systems in the literature however
have a more general elimination rule~\cite{DBLP:journals/pacmpl/AbelB20,grtt},
of the form:
\begin{align*}
\SYSTEMdruleGradBoxletGen{}
\end{align*}
%
The generalised rule allows an additional scaling of the resources
for $\SYSTEMnt{t_{{\mathrm{1}}}}$ due to additional use of $x$ in $\SYSTEMnt{t_{{\mathrm{2}}}}$. Such
a rule has not appeared in any Linear Base system, but could
easily be added to both systems to provide equivalent power
across both flavours of coeffect types.

This generalised typing enables the
derivation of a function $  \textcolor{coeffectColor}{\square_{ \SYSTEMnt{r} } }   (   \textcolor{coeffectColor}{\square_{ \SYSTEMnt{s} } }  \SYSTEMnt{A}   )    \multimap   \textcolor{coeffectColor}{\square_{ \SYSTEMsym{(}   \SYSTEMnt{r}  \cdot  \SYSTEMnt{s}   \SYSTEMsym{)} } }  \SYSTEMnt{A}  $,
akin to the multiplication (join) of a graded monad (the dual to comultiplication
of a graded comonad structure):
%
\begin{gather*}
  \vspace{-1em}
  \inferrule*[right=\SYSTEMRenameRuleGradBoxletGen{}]
             {\inferrule*[right=\SYSTEMRenameRuleGradvar{}]
                {\;}{   \SYSTEMmv{x}  :_{\textcolor{coeffectColor}{ \SYSTEMsym{1} } }    \textcolor{coeffectColor}{\square_{ \SYSTEMnt{r} } }   (   \textcolor{coeffectColor}{\square_{ \SYSTEMnt{s} } }  \SYSTEMnt{A}   )       \vdash_{\textsc{g}_\square }  \SYSTEMmv{x}  :   \textcolor{coeffectColor}{\square_{ \SYSTEMnt{r} } }   (   \textcolor{coeffectColor}{\square_{ \SYSTEMnt{s} } }  \SYSTEMnt{A}   )   }
               \inferrule*[right=\SYSTEMRenameRuleGradBoxletGen{}]
                          {\inferrule*[right=\SYSTEMRenameRuleGradvar{}]{\;}
                            {   \SYSTEMmv{y}  :_{\textcolor{coeffectColor}{ \SYSTEMsym{1} } }    \textcolor{coeffectColor}{\square_{ \SYSTEMnt{s} } }  \SYSTEMnt{A}      \vdash_{\textsc{g}_\square }  \SYSTEMmv{y}  :   \textcolor{coeffectColor}{\square_{ \SYSTEMnt{s} } }  \SYSTEMnt{A}  }
                            \inferrule*[right=\SYSTEMRenameRuleGradBoxpr{}]
                                       {\inferrule*[right=\SYSTEMRenameRuleGradvar{}]{\;}
                                         {   \SYSTEMmv{z}  :_{\textcolor{coeffectColor}{ \SYSTEMsym{1} } }  \SYSTEMnt{A}    \vdash_{\textsc{g}_\square }  \SYSTEMmv{z}  :  \SYSTEMnt{A} }}
                                       {   \SYSTEMmv{z}  :_{\textcolor{coeffectColor}{   \SYSTEMnt{r}  \cdot  \SYSTEMnt{s}   } }  \SYSTEMnt{A}    \vdash_{\textsc{g}_\square }   \textcolor{coeffectColor}{[}  \SYSTEMmv{z}  \textcolor{coeffectColor}{]}   :   \textcolor{coeffectColor}{\square_{   \SYSTEMnt{r}  \cdot  \SYSTEMnt{s}   } }  \SYSTEMnt{A}  }
                          \hspace{-2.4em}}
                          {   \SYSTEMmv{y}  :_{\textcolor{coeffectColor}{ \SYSTEMnt{r} } }    \textcolor{coeffectColor}{\square_{ \SYSTEMnt{s} } }  \SYSTEMnt{A}      \vdash_{\textsc{g}_\square }   \mathsf{let} \, \textcolor{coeffectColor}{[}  \SYSTEMmv{z}  \textcolor{coeffectColor}{]} =  \SYSTEMmv{y}  \, \mathsf{in} \,   \textcolor{coeffectColor}{[}  \SYSTEMmv{z}  \textcolor{coeffectColor}{]}    :   \textcolor{coeffectColor}{\square_{   \SYSTEMnt{r}  \cdot  \SYSTEMnt{s}   } }  \SYSTEMnt{A}  \hspace{-3em}}
             \hspace{-3em}}
             {    \SYSTEMmv{x}  :_{\textcolor{coeffectColor}{ \SYSTEMsym{1} } }    \textcolor{coeffectColor}{\square_{ \SYSTEMnt{r} } }   (   \textcolor{coeffectColor}{\square_{ \SYSTEMnt{s} } }  \SYSTEMnt{A}   )       \vdash_{\textsc{g}_\square }   \mathsf{let} \, \textcolor{coeffectColor}{[}  \SYSTEMmv{y}  \textcolor{coeffectColor}{]} =  \SYSTEMmv{x}  \, \mathsf{in} \,   \mathsf{let} \, \textcolor{coeffectColor}{[}  \SYSTEMmv{z}  \textcolor{coeffectColor}{]} =  \SYSTEMmv{y}  \, \mathsf{in} \,   \textcolor{coeffectColor}{[}  \SYSTEMmv{z}  \textcolor{coeffectColor}{]}     :   \textcolor{coeffectColor}{\square_{   \SYSTEMnt{r}  \cdot  \SYSTEMnt{s}   } }  \SYSTEMnt{A}   }
\end{gather*}
%
Indeed, the generalised $\Box$-elimination can be modelled in a
categorical semantics by a graded comonad along with an additional
graded multiplication operation $\mu_{r,s,A} :   \textcolor{coeffectColor}{\square_{ \SYSTEMnt{r} } }   (   \textcolor{coeffectColor}{\square_{ \SYSTEMnt{s} } }  \SYSTEMnt{A}   )    \multimap   \textcolor{coeffectColor}{\square_{ \SYSTEMsym{(}   \SYSTEMnt{r}  \cdot  \SYSTEMnt{s}   \SYSTEMsym{)} } }  \SYSTEMnt{A}  $ (as one might find in a graded monad structure).
However, not all graded comonads modelling $ \textcolor{coeffectColor}{\square_{ \SYSTEMnt{r} } }  \SYSTEMnt{A} $
admit such an operation and so having the generalised
$\Box$-elimination may not be desirable in general graded comonadic calculi.
We thus omitted this generalised rule from the main part of our results, considering
the simpler $\SYSTEMRenameRuleLinlet{}$ and $\SYSTEMRenameRuleGradBoxlet{}$ rules.

